\section{Five open questions}

\begin{enumerate}
  \item \textbf{Non-abelian generalisation.}  Does the exponential-in-$D$ convergence
    and shallow $D$-vs-$N$ slope survive extension to non-abelian gauge groups
    (SU(2), SU(3)) where link variables force a larger local Hilbert space and
    Gauss-law constraints become non-trivial?  \emph{Basis:} the paper treats
    only U(1) 1+1D (Schwinger); the larger $H_0$ spectrum could both reduce
    reference-state / ground-state overlap and worsen Hankel ill-conditioning.
    \emph{Next steps:} pick an SU(2) Kogut--Susskind lattice ($N=4$, $j_{\max}=1/2$
    truncation), build $H$ analogous to Eq.~15, choose a strong-coupling reference
    state, re-run the QSE-vs-$D$ sweep classically, and compare threshold-$D$ to
    the U(1) case at matched qubit count.

  \item \textbf{Comparison to modern quantum-Krylov variants.}  How does the paper's
    Hankel/moment form compare quantitatively to QLanczos (Motta et al.),
    real-time-evolution QK-eig (Klymko et al.), and thermal Krylov variants on
    the same Eq.~15 instance?  \emph{Basis:} the paper uses one specific quantum
    Krylov protocol; other variants trade different resources (imaginary time,
    Trotterised real time, thermal states) and their relative advantage on this
    Schwinger instance is untested.  \emph{Next steps:} implement each variant
    classically on Eq.~15 at $\mu = 1.5$, $x = 0.5$, $N = 4$--$10$; compare
    threshold-$D$, $\kappa(S)$ growth, and shot-noise sensitivity.

  \item \textbf{Quantitative resource-estimate verification.}  Is the block-encoding
    $\tilde O(N)$ gate-cost advantage over prior $\tilde O(N^2)$ LCU quantitatively
    confirmed when the qubitization circuit is compiled and T-counted for the
    paper's own target sizes?  \emph{Basis:} Sec.~5--6 gives an asymptotic
    analysis; constants and log factors matter for crossover.  \emph{Next steps:}
    use Qualtran or Azure QRE to compile the PREPARE/SELECT decomposition for
    Eq.~15 at $N = 10, 20, 50, 100$; T-count the block encoding; compare against
    naive-LCU compilation; report the crossover $N$.

  \item \textbf{Noise robustness of the Krylov step.}  How does the algorithm
    degrade under realistic gate noise (not just shot noise on moments), and at
    what noise level does exponential-in-$D$ convergence collapse?  \emph{Basis:}
    Sec.~4.3 shot-noise extrapolation does not model per-gate depolarising or
    dephasing; block-encoding errors compound with moment order $k$ and can cap
    the effective usable $D$ independently of shot count.  \emph{Next steps:}
    simulate noisy moment estimation on Eq.~15 with shot noise plus per-gate
    depolarising rates $10^{-3}, 10^{-4}, 10^{-5}$; plot smallest achievable
    $\Delta E / E_{\mathrm{int}}$ vs noise rate; estimate $D_{\max}(\mathrm{noise})$
    and compare to $D(N) \approx 0.057 N + 4.36$.

  \item \textbf{ML-driven Krylov subspace selection.}  Can a learned / adaptive
    basis choice (restarted Krylov, block Krylov, or a variational reference-state
    ansatz outer-optimised before the Krylov step) reduce the required $D$ or
    delay Hankel breakdown for the same instance?  \emph{Basis:} in classical
    numerical linear algebra, restarted and block Krylov methods routinely
    reduce $D$ substantially at fixed accuracy; a quantum-implementable analogue
    of these ideas for QSE is essentially unexplored.  \emph{Next steps:}
    classically explore restarted Krylov (keeping $K$ best Ritz vectors), block
    Krylov (multiple orthogonal reference states), and a small parametric
    reference-state ansatz optimised by a VQE outer loop; compare threshold-$D$
    and $\kappa(S)$ growth to the paper's fixed-basis Hankel baseline.
\end{enumerate}
